\renewcommand{\figurename}{Abb.}

\section{Einleitung}
\label{sec:einleitung}

\begin{quote}
\textit{%
The Covid-19 pandemic has caused more disruption to the energy sector than any other event in recent history, leaving impacts that will be felt for years to come.
}\\
-- IEA 2020~\cite{iea2020}
\end{quote}

\subsection{Problemstellung und Zielsetzung}
\label{subsec:problemstellung}

Die COVID-19 Krise hat nahezu alle Bereiche unserer modernen Gesellschaft auf die Probe gestellt, besonders aber den Energiesektor. Dies hat die \textbf{International Energy Agency} in ihrem Bericht \textit{World Energy Outlook} von 2020 festgestellt. Dieser befindet sich in einem tiefgreifenden Wandel, geprägt von Dekarbonisierung, Digitalisierung und geopolitischen Unsicherheiten. Prognosemodelle für Energieverbrauch sind ein zentrales Werkzeug der, um Entwicklungen oder Trends frühzeitig zu erkennen. Die Anwendung fortgeschrittener Zeitreihenmodelle trägt zur strategischen Planung im Kontext langfristiger Trends und plötzlich eintretender externer Schocks bei. Der methodische Vergleich von Prognoseverfahren und die Fehleranalyse bei außergewöhnlichen Ereignissen sind daher von hoher praktischer und wissenschaftlicher Relevanz.


\subsection{Aufbau der Arbeit}
\label{subsec:aufbau}

Diese Arbeit befasst sich mit der Frage, ob Zeitreihenanalysen in der Lage sind, ihre Vorhersagen auf die Schocks externer Ereignisse anzupassen. Mittels der Analyse und Prognose des US-amerikanischen Primärenergieverbrauchs auf Basis von ARIMA-Modellen und dem Holt-Winters Algorithmus werden die Auswirkungen globaler Krisen wie COVID-19 oder der russische Invasionskrieg in der Ukraine auf die Genauigkeit mathematischer Vorhersagen betrachtet. Ziel ist es, die verschiedenen Modellierungsansätze zu vergleichen und deren Prognosegüte systematisch zu bewerten.

Die Arbeit ist wie folgt gegliedert: Zunächst werden die theoretischen und methodischen Grundlagen im Bereich Zeitreihenanalyse und Prognosemetriken erläutert. Im Anschluss folgt eine detaillierte Beschreibung des verwendeten Datensatze sowie der Implementierung der verwendeten Prognosemodelle und deren Validierung vor. Im Ergebnisteil werden die Prognosegüte und die Fehlerverteilung analysiert und durch verschiedene Visualisierungen veranschaulicht. Anschließend erfolgt eine vertiefte Betrachtung der Prognosefehler während relevanter externer Ereignisse.
